\newcommand{\ModuleID}{EL-I.5}
\newcommand{\ModuleTitle}{Fourth and Fifth Declensions}
\newcommand{\ModuleStage}{Stage I — Foundations}
\newcommand{\ModuleUnit}{I.5}
\newcommand{\ModuleOutcome}{Recognize fourth- and fifth-declension forms and discriminate look-alike endings across all five noun declensions}
\newcommand{\ModulePrerequisites}{EL-I.4 mastery: recover third-declension stems and apply neuter and i-stem rules}
\newcommand{\ModuleSessionOne}{Learn the masculine and neuter fourth-declension and regular fifth-declension patterns, then complete Worksheet I.17 as the guided application.}
\newcommand{\ModuleSessionTwo}{Retrieve the five-declension look-alike ending matrix, use syntax to resolve ambiguous forms, complete Worksheet I.20 independently, and document each decision.}
\newcommand{\ModulePrintedPractice}{Worksheets I.17 (guided Variant A) and I.20 (independent Variant D), with terminal keys}
\newcommand{\ModuleExtensionPractice}{Worksheets I.18 and I.19 remain optional remediation or delayed-recall variants}
\newcommand{\ModuleNext}{EL-I.6, Adjectives and Agreement}
\newcommand{\ModuleMastery}{Write the three displayed paradigms with at least 95 percent accuracy, list all possible analyses for 18 of 20 mixed-declension forms, and select the contextual analysis in at least nine of ten short phrases.}
\newcommand{\ModuleDelayNote}{}
\newcommand{\ModuleReferenceFocus}{%
  \begin{itemize}
    \item fourth-declension masculine and neuter dictionary forms and paradigms;
    \item the fifth-declension pattern represented by \Latin{res, rei};
    \item high-risk \Latin{-us}, \Latin{-u}, \Latin{-ui}, \Latin{-us} genitive, and \Latin{-es/-ei} forms; and
    \item a five-declension workflow: dictionary genitive, possible forms, syntax, agreement, and government.
  \end{itemize}}
\newcommand{\ModuleContrast}{A familiar-looking nominative does not identify a declension. The dictionary genitive establishes the class; syntax then chooses among forms that share an ending.}
\newcommand{\ModuleLessonSource}{\CourseRoot/shared/content/volume1/all}
\newcommand{\ModuleExerciseSource}{\CourseRoot/shared/exercises/volume1/all}
\newcommand{\ModuleWorkedBank}{I}
\newcommand{\ModuleExerciseBank}{I}
\newcommand{\ModuleWorksheetVariants}{A,D}
\newcommand{\ModuleScopeNote}{This packet completes the ordinary five-declension map while keeping rare lexical behavior dependent on the identified dictionary entry.}
\newcommand{\ModuleEnrichment}{%
  \SubmoduleExtension
    {Fourth-declension forms}
    {Fourth-declension \Latin{-us} and \Latin{-u} forms are best learned as coordinates in a paradigm, not as isolated endings.}
    {Make a two-column matrix for the lesson's masculine and neuter fourth-declension models. Highlight every nominative--accusative identity and every form that can express more than one case or number.}
    {State the neuter rule and identify the dictionary form that distinguishes fourth declension from a superficially similar second-declension noun.}
    {The matrix must preserve the lesson paradigms and mark neuter nominative and accusative identity. The genitive singular in \Latin{-us}, together with the full dictionary pair, identifies fourth declension; the nominative alone does not.}
  \SubmoduleExtension
    {Fifth-declension forms}
    {The small fifth declension contains high-frequency words whose forms can resemble endings from other classes.}
    {Build the full lesson paradigm for \Latin{res} from memory, then annotate which cells would remain ambiguous if the dictionary entry and clause were hidden.}
    {Give the genitive and dative singular endings and explain how agreement can help distinguish a fifth-declension form in context.}
    {The paradigm must give genitive and dative singular \Latin{rei}. A correct explanation uses an agreeing adjective or demonstrative plus clause function to narrow the analysis; it does not claim that one English gloss proves the case.}
  \SubmoduleExtension
    {A five-declension decision procedure}
    {A look-alike ending matrix makes ambiguity explicit before context resolves it.}
    {Create rows for \Latin{-a, -æ, -i, -is, -us, -u, -es}, and enter at least two possible declension/case/number analyses where the lesson permits them. Apply the matrix to the lesson's five-declension review phrases.}
    {Recite the decision order: dictionary evidence, morphological possibilities, clause function, agreement, and government.}
    {The matrix must not collapse genuinely different analyses. Each applied decision must cite the noun's dictionary genitive first and then syntactic evidence. The stated order must preserve morphology before translation.}
  \SubmoduleExtension
    {Handwritten study and controlled composition}
    {Number transformations expose whether the learner controls both rare paradigms and their agreement partners.}
    {Complete the lesson's \Latin{dies} and \Latin{genu} transformations, drawing a box around every noun ending and an arrow to each word whose agreement must also change.}
    {Name one unchanged lexical stem and one changed grammatical coordinate in each transformation.}
    {The lesson models yield plural \Latin{Hi sunt dies} and \Latin{omnia genua}. The analysis must keep the stems stable, identify number as the changed coordinate, and account for agreement in the accompanying demonstrative, verb, or adjective.}
}
